% SGA 3, Expose VI_B, Loop 1 English reconstruction.
% Scope: paragraph 11.4, Proposition 11.5, paragraph 11.6,
% Remark 11.6.A, Proposition 11.6.1, Lemma 11.7, and
% Remark 11.7.bis, including editor's notes 108--111.
% Authority: Polo--Gille Exp6B-13oct24.pdf, component-local
% pages 78--82 (combined-reader pages 419--423; re-edition
% headers 408--412; original-margin printed pages 397--401).
% Paragraph 11.4 begins on local page 78, still within printed page 397;
% printed page 398 begins there only at Proposition 11.5. Definition
% 11.8.0 begins on local page 82, still within printed page 401, and is
% excluded. Printed page 402 begins on local page 85.
% OCR and Reinhold English Markdown are comparison witnesses only.
% Loop 2 replaces all seven source-bound commutative-diagram crops
% with faithful native LaTeX reconstructions; all are closed.

\medskip
\phantomsection\label{par:sga3-VIB-11.4}%
\noindent\textbf{11.4.---}%
\footnote{\label{note:sga3-VIB-108}编者注: 在\SGARefLink{par:sga3-VIB-11.4}{11.4}--\SGARefLink{par:sga3-VIB-11.6}{11.6} 中,
我们考虑的是
\(\underline{\operatorname{Hom}}_{\mathcal O_S}
  (\mathrm W(\mathcal E),\mathrm W(\mathcal F))\)
而不是
\(\underline{\operatorname{Hom}}_{\mathcal O_S}
  (\mathrm V(\mathcal F),\mathrm V(\mathcal E))\)
\textup{(参见 I, \SGARefLink{sga3:i:proposition:4-6-3}{4.6.3})},
并利用增补的\hyperref[concl:sga3-VIB-11.0-b]{11.0}\textup{(b)} 简化了原版.}
设 \(\mathcal E\) 和 \(\mathcal F\) 为两个拟凝聚 \(\mathcal O_S\)-模.
考虑 \(S\)-函子
\[
  \underline{\operatorname{Hom}}_{\mathcal O_S}
    \bigl(\mathrm W(\mathcal E),\mathrm W(\mathcal F)\bigr)
\]
\textup{(参见 I, \SGARefLink{sga3:i:definition:3-1-4}{3.1.4})},
即对每个 \(S\)-概形 \(f:X\to S\),
\phantomsection\label{form:sga3-VIB-11.4-HomX}%
\[
  \underline{\operatorname{Hom}}_{\mathcal O_S}
    \bigl(\mathrm W(\mathcal E),\mathrm W(\mathcal F)\bigr)(X)
  =
  \operatorname{Hom}_{\mathcal O_X}
    \bigl(\mathrm W(\mathcal E)_X,\mathrm W(\mathcal F)_X\bigr).
\]
此外, 由 I, \SGARefLink{sga3:i:proposition:4-6-2}{4.6.2}, 有
\(\mathrm W(\mathcal E)_X=\mathrm W(f^*(\mathcal E))\)
\textup{(对 \(\mathcal F\) 同样如此)}, 且
\[
  \operatorname{Hom}_{\mathcal O_X}
    \bigl(\mathrm W(f^*(\mathcal E)),
          \mathrm W(f^*(\mathcal F))\bigr)
  =
  \operatorname{Hom}_{\mathcal O_X}
    \bigl(f^*(\mathcal E),f^*(\mathcal F)\bigr).
\]
因此得到 \textup{(最后一个等式使用伴随公式)}
\phantomsection\label{form:sga3-VIB-11.4-adjunction}%
\[
  \tag{\(\dagger\)}
  \underline{\operatorname{Hom}}_{\mathcal O_S}
    \bigl(\mathrm W(\mathcal E),\mathrm W(\mathcal F)\bigr)(X)
  =
  \operatorname{Hom}_{\mathcal O_X}
    \bigl(f^*(\mathcal E),f^*(\mathcal F)\bigr)
  =
  \operatorname{Hom}_{\mathcal O_S}
    \bigl(\mathcal E,f_*f^*(\mathcal F)\bigr).
\]

\medskip
\phantomsection\label{prop:sga3-VIB-11.5}%
\noindent\textbf{命题 11.5.---}
\emph{设 \(X\) 为在 \(S\) 上拟紧且拟分离的 \(S\)-概形,
\(f:X\to S\) 为结构态射,
\(\mathcal E\) 和 \(\mathcal E'\) 为两个拟凝聚 \(\mathcal O_S\)-模.
假设下列两个条件之一成立:}
\begin{enumerate}
\item[\textup{a)}]\phantomsection\label{concl:sga3-VIB-11.5-a}
  \emph{\(f\) 仿射;}

\item[\textup{b)}]\phantomsection\label{concl:sga3-VIB-11.5-b}
  \emph{\(\mathcal E\) 在 \(\mathcal O_S\) 上平坦.}
\end{enumerate}
\emph{则典范态射}
\phantomsection\label{form:sga3-VIB-11.5-projection}%
\[
  \mathcal E\otimes_{\mathcal O_S}\mathcal A(X)
    \longrightarrow f_*f^*(\mathcal E)
\]
\emph{为同构, 从而}
\phantomsection\label{form:sga3-VIB-11.5-Hom}%
\[
  \underline{\operatorname{Hom}}_{\mathcal O_S}
    \bigl(\mathrm W(\mathcal E'),\mathrm W(\mathcal E)\bigr)(X)
  =
  \operatorname{Hom}_{\mathcal O_S}
    \bigl(\mathcal E',
          \mathcal E\otimes_{\mathcal O_S}\mathcal A(X)\bigr).
\]

事实上, 第二个断言由\hyperref[par:sga3-VIB-11.4]{11.4} 和第一个断言得到;
而第一个断言在情形~\hyperref[concl:sga3-VIB-11.5-a]{\textup{(a)}} 中
由 EGA~II, 1.4.7 得到,
在情形~\hyperref[concl:sga3-VIB-11.5-b]{\textup{(b)}} 中
由\hyperref[concl:sga3-VIB-11.0-b]{11.0\textup{(b)}} 得到.% 
\footnote{\label{note:sga3-VIB-109}编者注: 我们简化了原版的论证; 原版先使用同构
\[
  \underline{\operatorname{Hom}}_{\mathcal O_S}
    \bigl(\mathrm W(\mathcal E'),\mathrm W(\mathcal E)\bigr)
  \simeq
  \underline{\operatorname{Hom}}_{\mathcal O_S}
    \bigl(\mathrm V(\mathcal E),\mathrm V(\mathcal E')\bigr),
\]
再利用右端到下列对象的包含映射:
\[
  \underline{\operatorname{Hom}}_S
    \bigl(\mathrm V(\mathcal E),\mathrm V(\mathcal E')\bigr)
  =
  \operatorname{Hom}_{\mathcal O_S}
    \bigl(\mathcal E',
          f_*f^*(\operatorname{Sym}(\mathcal E))\bigr),
\]
并对 \(\mathrm V(\mathcal E)=
  \operatorname{Spec}(\operatorname{Sym}(\mathcal E))\)
应用 EGA~III, 4.1.15, 从而推出
\SGARefLink{par:sga3-VIB-11.0}{11.0}\textup{\SGARefLink{par:sga3-VIB-11.0:item:b}{(b)}}.}

\medskip
\phantomsection\label{par:sga3-VIB-11.6}%
\noindent\textbf{11.6.---}%
\footnote{\label{note:sga3-VIB-110}编者注: 扩充了\SGARefLink{par:sga3-VIB-11.6}{11.6},
并将所得结果写成\hyperref[prop:sga3-VIB-11.6.1]{命题 11.6.1} 的形式.}
设 \(G\) 为 \(S\)-群, \(\mathcal E\) 为拟凝聚 \(\mathcal O_S\)-模.
给 \(\mathcal E\) 赋予 \(G\)-\(\mathcal O_S\)-模结构
\textup{(即 \(G\) 在 \(\mathrm W(\mathcal E)\) 上的 \(\mathcal O_S\)-线性作用,
参见 I, \SGARefLink{sga3:i:definition:4-7-1}{4.7.1})},
等价于给定幺半群值 \(S\)-函子的态射
\phantomsection\label{form:sga3-VIB-11.6-rho}%
\[
  \rho:G\longrightarrow
  \underline{\operatorname{End}}_{\mathcal O_S}
    \bigl(\mathrm W(\mathcal E)\bigr);
\]
事实上, 这样的 \(\rho\) 必将 \(G\) 映入
\(\underline{\operatorname{Aut}}_{\mathcal O_S}
  (\mathrm W(\mathcal E))\).

由\hyperref[par:sga3-VIB-11.4]{11.4}, 给定 \(S\)-函子的态射
\[
  \rho:G\longrightarrow
  \underline{\operatorname{End}}_{\mathcal O_S}
    \bigl(\mathrm W(\mathcal E)\bigr)
\]
等价于给定下列集合的元素 \(\theta\):
\[
  \operatorname{Hom}_{\mathcal O_G}
    \bigl(f^*(\mathcal E),f^*(\mathcal E)\bigr),
\]
由伴随性, 此元素对应于下列集合的元素 \(\delta\):
\[
  \operatorname{Hom}_{\mathcal O_S}
    \bigl(\mathcal E,f_*f^*(\mathcal E)\bigr),
\]
其中 \(f\) 表示投影 \(G\to S\).

设 \(m:G\times_SG\to G\) 为乘法, \(\delta_G\) 为态射
\[
  \mathcal O_G\longrightarrow
  m_*(\mathcal O_{G\times_SG}),
\]
并设 \(\varphi:G\times_SG\to S\) 为投影
\textup{(它等于 \(f\circ m\))}.
为方便起见, 用 \(\boxtimes\) 表示``外''张量积. 这样得到态射
\phantomsection\label{form:sga3-VIB-11.6-id-deltaG}%
\[
  \operatorname{id}_{\mathcal E}\boxtimes\delta_G:
  f^*(\mathcal E)
    =\mathcal E\boxtimes_{\mathcal O_S}\mathcal O_G
  \longrightarrow
  \mathcal E\boxtimes_{\mathcal O_S}
    m_*(\mathcal O_{G\times_SG}).
\]
滥用记号, 仍用 \(\operatorname{id}_{\mathcal E}\boxtimes\delta_G\)
表示此态射与典范态射
\[
  \mathcal E\boxtimes_{\mathcal O_S}
    m_*(\mathcal O_{G\times_SG})
  \longrightarrow
  m_*m^*f^*(\mathcal E)
  =m_*\varphi^*(\mathcal E).
\]

的复合. 另一方面, 用 \(h:G\to S\) 表示 \(f:G\to S\) 的第二个副本,
并考虑下述交换图, 其中 \(p,q\) 表示两个投影:
\begin{center}
  \phantomsection\label{diag:sga3-VIB-11.6-projections}%
  \begin{tikzcd}[row sep=large,column sep=large]
    G\times_SG
      \arrow[r,"q"]
      \arrow[d,"p"']
      \arrow[dr,"\varphi" description] &
    G \arrow[d,"h"] \\
    G \arrow[r,"f"'] &
    S
  \end{tikzcd}
\end{center}
仍用 \(\delta\) 表示态射 \(\mathcal E\to h_*h^*(\mathcal E)\), 得到态射
\phantomsection\label{form:sga3-VIB-11.6-delta-id}%
\[
  \delta\boxtimes\operatorname{id}_{\mathcal O_G}:
  f^*(\mathcal E)
    =\mathcal E\boxtimes_{\mathcal O_S}\mathcal O_G
  \longrightarrow
  h_*h^*(\mathcal E)\boxtimes_{\mathcal O_S}\mathcal O_G
    =f^*h_*h^*(\mathcal E).
\]
滥用记号, 仍用 \(\delta\boxtimes\operatorname{id}_{\mathcal O_G}\)
表示此态射与典范态射
\[
  f^*h_*h^*(\mathcal E)\longrightarrow
  p_*\varphi^*(\mathcal E).
\]

的复合. 于是 \(\rho\) 与乘法相容这一条件,
等价于对 \(S\) 的每个开子集 \(U\), 下图交换:
\begin{center}
  \phantomsection\label{diag:sga3-VIB-11.6-1}%
  \makebox[0pt][r]{\textup{(1)}\qquad}%
  {\scriptsize
  \begin{tikzcd}[row sep=large,column sep=large]
    \Gamma(U,\mathcal E)
      \arrow[r,"\delta"]
      \arrow[d,"\delta"'] &
    \Gamma\!\left(
      U\times_SG,\,
      \mathcal E\boxtimes_{\mathcal O_S}\mathcal O_G
    \right)
      \arrow[d,"{\operatorname{id}_{\mathcal E}\boxtimes\delta_G}"] \\
    \Gamma\!\left(
      U\times_SG,\,
      \mathcal E\boxtimes_{\mathcal O_S}\mathcal O_G
    \right)
      \arrow[r,"{\delta\boxtimes\operatorname{id}_{\mathcal O_G}}"'] &
    \Gamma\!\left(
      U\times_SG\times_SG,\,
      \mathcal E\boxtimes_{\mathcal O_S}\mathcal O_G
        \boxtimes_{\mathcal O_S}\mathcal O_G
    \right)
  \end{tikzcd}}
\end{center}
此外, 单位截面 \(\varepsilon:S\to G\)
诱导从 \(\mathcal O_G\) 到下列对象的态射 \(u\):
\[
  \varepsilon_*\varepsilon^*(\mathcal O_G)
  =\varepsilon_*(\mathcal O_S),
\]
而 \(\rho\) 保持单位元这一条件等价于下图交换:
\begin{center}
  \phantomsection\label{diag:sga3-VIB-11.6-2}%
  \makebox[0pt][r]{\textup{(2)}\qquad}%
  {\small
  \begin{tikzcd}[row sep=large,column sep=large]
    \Gamma(U,\mathcal E)
      \arrow[rr,"\delta"]
      \arrow[dr,"\sim"'] &
    {} &
    \Gamma\!\left(
      U\times_SG,\,
      \mathcal E\boxtimes_{\mathcal O_S}\mathcal O_G
    \right)
      \arrow[dl,"{\operatorname{id}_{\mathcal E}\boxtimes u}"] \\
    {} &
    \Gamma\!\left(
      U\times_SS,\,
      \mathcal E\boxtimes_{\mathcal O_S}\mathcal O_S
    \right)\mathpunct{.} &
    {}
  \end{tikzcd}}
\end{center}

因此, 给 \(\mathcal E\) 赋予 \(G\)-\(\mathcal O_S\)-模结构,
等价于给定 \(\mathcal O_S\)-模态射
\[
  \delta:\mathcal E\longrightarrow f_*f^*(\mathcal E)
\]
满足上面的条件\hyperref[diag:sga3-VIB-11.6-1]{\textup{(1)}}
和\hyperref[diag:sga3-VIB-11.6-2]{\textup{(2)}}. 此时由 \(\delta\)
经伴随性得到的态射
\[
  \theta:f^*(\mathcal E)\longrightarrow f^*(\mathcal E)
\]
为同构, 因为它对应于下述同构:
\[
  G\times_S\mathrm W(\mathcal E)
    \xrightarrow{\ \sim\ }
  G\times_S\mathrm W(\mathcal E)
\]
此同构在集合意义下由 \((g,x)\longmapsto(g,gx)\) 定义;
亦见 I, \SGARefLink{sga3:i:remark:6-5-4}{6.5.4}.

现在假设 \(G\) 在 \(S\) 上平坦、拟紧且拟分离,
并假设 \(\mathcal A(G)\) 为平坦 \(\mathcal O_S\)-模.
则由\hyperref[concl:sga3-VIB-11.1-c]{11.1\textup{(c)}}, 典范态射
\[
  \mathcal A(G)\otimes_{\mathcal O_S}\mathcal A(G)
    \longrightarrow
  \mathcal A(G\times_SG)
\]
为同构, 用 \(\Delta\) 表示态射
\[
  \delta_G:\mathcal A(G)\longrightarrow
  \mathcal A(G)\otimes_{\mathcal O_S}\mathcal A(G).
\]

此外, 若
\[
  \mathcal E\otimes_{\mathcal O_S}\mathcal A(G)
    =f_*f^*(\mathcal E),
\]
如\hyperref[prop:sga3-VIB-11.5]{11.5} 所示,
当 \(G\to S\) 仿射或 \(\mathcal E\) 在 \(\mathcal O_S\) 上平坦时即如此,
则条件\hyperref[diag:sga3-VIB-11.6-1]{\textup{(1)}}
和\hyperref[diag:sga3-VIB-11.6-2]{\textup{(2)}} 等价于下述条件,
它们表示
\[
  \delta:\mathcal E\longrightarrow
  \mathcal E\otimes_{\mathcal O_S}\mathcal A(G)
\]
将 \(\mathcal E\) 赋成右 \(\mathcal A(G)\)-余模
\textup{(参见 I, \SGARefLink{sga3:i:proposition:4-7-2}{4.7.2})}:

\smallskip
\phantomsection\label{cond:sga3-VIB-11.6-CM1}%
\noindent\textup{(CM 1)}
令 \(\mathcal A=\mathcal A(G)\). 下图交换:
\begin{center}
  \phantomsection\label{diag:sga3-VIB-11.6-CM1}%
  \begin{tikzcd}[row sep=large,column sep=large]
    \mathcal E
      \arrow[r,"\delta"]
      \arrow[d,"\delta"'] &
    \mathcal E\otimes\mathcal A
      \arrow[d,"{\operatorname{id}_{\mathcal E}\otimes\Delta}"] \\
    \mathcal E\otimes\mathcal A
      \arrow[r,"{\delta\otimes\operatorname{id}_{\mathcal A}}"'] &
    \mathcal E\otimes\mathcal A\otimes\mathcal A
  \end{tikzcd}
\end{center}

\smallskip
\phantomsection\label{cond:sga3-VIB-11.6-CM2}%
\noindent\textup{(CM 2)}
用 \(\eta:\mathcal A\to\mathcal O_S\) 表示态射
\(\mathcal A(\varepsilon)\). 下图交换:
\begin{center}
  \phantomsection\label{diag:sga3-VIB-11.6-CM2}%
  \begin{tikzcd}[row sep=large,column sep=large]
    \mathcal E
      \arrow[rr,"\delta"]
      \arrow[dr,"\sim"'] &
    {} &
    \mathcal E\otimes\mathcal A
      \arrow[dl,"{\operatorname{id}_{\mathcal E}\otimes\eta}"] \\
    {} &
    \mathcal E\otimes\mathcal O_S\mathpunct{.} &
    {}
  \end{tikzcd}
\end{center}

\medskip
\phantomsection\label{rem:sga3-VIB-11.6.A}%
\noindent\textbf{注 11.6.A.---}
回顾, \(\mathrm V(\mathcal E)\) 表示表出函子
\(\mathrm V(\mathcal E)\) 的 \(S\) 上的向量纤维化;
即对每个 \(S'\to S\),
\[
  \mathrm V(\mathcal E)(S')
  =
  \operatorname{Hom}_{\mathcal O_{S'}}
    \bigl(\mathcal E_{S'},\mathcal O_{S'}\bigr).
\]
由于 I, \SGARefLink{sga3:i:proposition:4-6-2}{4.6.2}
给出幺半群值 \(S\)-函子的反同构
\[
  \underline{\operatorname{End}}_{\mathcal O_S}
    \bigl(\mathrm W(\mathcal E)\bigr)
  \simeq
  \underline{\operatorname{End}}_{\mathcal O_S}
    \bigl(\mathrm V(\mathcal E)\bigr),
\]
若 \(\mathcal E\) 为左 \(G\)-\(\mathcal O_S\)-模,
则 \(G\) 在 \(\mathrm V(\mathcal E)\) 上有右作用
\[
  \mu:\mathrm V(\mathcal E)\times_SG
    \longrightarrow\mathrm V(\mathcal E)
\]
在集合意义下定义为
\[
  (\phi g)(x)=\phi(gx)
\]
其中 \(g\in G(S')\), \(x\in\Gamma(S',\mathcal E_{S'})\),
\(\phi\in
  \operatorname{Hom}_{\mathcal O_{S'}}
    (\mathcal E_{S'},\mathcal O_{S'})\) 任意.
因此有下述交换图
\begin{center}
  \phantomsection\label{diag:sga3-VIB-11.6.A-actions}%
  {\small
  \begin{tabular}{c@{\qquad}c}
    \begin{tikzcd}[row sep=large,column sep=large]
      \mathrm V(\mathcal E)\times_SG\times_SG
        \arrow[r,"{\mu\times\operatorname{id}_G}"]
        \arrow[d,"{\operatorname{id}_{\mathrm V(\mathcal E)}\times m}"'] &
      \mathrm V(\mathcal E)\times_SG
        \arrow[d,"\mu"] \\
      \mathrm V(\mathcal E)\times_SG
        \arrow[r,"\mu"'] &
      \mathrm V(\mathcal E)
    \end{tikzcd}
    &
    \begin{tikzcd}[row sep=large,column sep=large]
      \mathrm V(\mathcal E) &
      \mathrm V(\mathcal E)\times_SG
        \arrow[l,"\mu"'] \\
      {} &
      \mathrm V(\mathcal E)\times_SS\mathpunct{.}
        \arrow[u,"{\operatorname{id}_{\mathrm V(\mathcal E)}\times\varepsilon}"]
        \arrow[ul,"\sim"']
    \end{tikzcd}
  \end{tabular}}
\end{center}
当 \(G\) 在 \(S\) 上平坦、拟紧且拟分离,
\(\mathcal A(G)\) 为平坦 \(\mathcal O_S\)-模,
且满足\hyperref[prop:sga3-VIB-11.5]{11.5} 的条件之一时,
同样可重新得到条件\hyperref[cond:sga3-VIB-11.6-CM1]{\textup{(CM 1)}}
和\hyperref[cond:sga3-VIB-11.6-CM2]{\textup{(CM 2)}}.

因此, 我们证明了:

\medskip
\phantomsection\label{prop:sga3-VIB-11.6.1}%
\noindent\textbf{命题 11.6.1.---}
\emph{设 \(G\) 为在 \(S\) 上平坦、拟紧且拟分离的 \(S\)-群,
使 \(\mathcal A(G)\) 为平坦 \(\mathcal O_S\)-模,
并设 \(\mathcal E\) 为拟凝聚 \(\mathcal O_S\)-模.}
\begin{enumerate}
\item[\textup{(i)}]\phantomsection\label{concl:sga3-VIB-11.6.1-i}
  \emph{给定 \(\mathcal A(G)\)-余模结构}
  \[
    \delta:\mathcal E\longrightarrow
      \mathcal E\otimes_{\mathcal O_S}\mathcal A(G)
  \]
  \emph{等价于给 \(\mathcal E\) 赋予 \(G_{\mathrm{af}}\)-\(\mathcal O_S\)-模结构
\textup{(即 \(G_{\mathrm{af}}\) 在 \(\mathcal E\) 上的 \(\mathcal O_S\)-线性作用)}.
与 \(S\)-群态射 \(G\to G_{\mathrm{af}}\) 复合,
便在 \(\mathcal E\) 上定义了 \(G\)-\(\mathcal O_S\)-模结构.}

\item[\textup{(ii)}]\phantomsection\label{concl:sga3-VIB-11.6.1-ii}
  \emph{此外, 若 \(\mathcal E\) 平坦, 则 \(G\) 在 \(\mathcal E\) 上的每个
\(\mathcal O_S\)-线性作用都经由 \(G_{\mathrm{af}}\) 分解,
并对应于 \(\mathcal E\) 上唯一的 \(\mathcal A(G)\)-余模结构.}
\end{enumerate}

\medskip
\phantomsection\label{lem:sga3-VIB-11.7}%
\noindent\textbf{引理 11.7.---}
\emph{设 \(G\) 为在 \(S\) 上平坦、拟紧且拟分离的 \(S\)-群,
使 \(\mathcal A=\mathcal A(G)\) 为平坦 \(\mathcal O_S\)-模.
设 \(\mathcal E\) 为拟凝聚 \(\mathcal O_S\)-模,}
\[
  \delta:\mathcal E\longrightarrow\mathcal E\otimes\mathcal A
\]
\emph{为 \(\mathcal A\)-余模结构, 且}
\[
  \rho:G\longrightarrow
  \underline{\operatorname{Aut}}_{\mathcal O_S}
    \bigl(\mathrm W(\mathcal E)\bigr)
\]
\emph{为关联的 \(G\) 在 \(\mathcal E\) 上的作用.
设 \(\mathcal E_0\) 为 \(\mathcal E\) 的拟凝聚 \(\mathcal O_S\)-子模,
使 \(\delta\) 在 \(\mathcal E_0\) 上的限制 \(\delta_0\)}%
\begingroup
\renewcommand{\thefootnote}{\fnsymbol{footnote}}%
\footnote[1]{\label{note:sga3-VIB-11.7-delta-source}
英译者注: 法文底本在此处作``\(\mu\) 的限制 \(\delta_0\)''.
根据前后的记号和所列方图, 此处应为 \(\delta\).}%
\endgroup
\emph{经由下述态射分解}
\[
  \mathcal E_0\longrightarrow\mathcal E_0\otimes\mathcal A;
\]
\emph{即假设下图交换:}
\begin{center}
  \phantomsection\label{diag:sga3-VIB-11.7-stability}%
  \begin{tikzcd}[row sep=large,column sep=large]
    \mathcal E_0
      \arrow[r,hook]
      \arrow[d,"\delta_0"'] &
    \mathcal E
      \arrow[d,"\delta"] \\
    \mathcal E_0\otimes\mathcal A
      \arrow[r,hook] &
    \mathcal E\otimes\mathcal A\mathpunct{.}
  \end{tikzcd}
\end{center}
\emph{\textup{(注意: 态射 \(\mathcal E_0\otimes\mathcal A
  \to\mathcal E\otimes\mathcal A\)
为单射, 因为 \(\mathcal A\) 在 \(\mathcal O_S\) 上平坦.)}
则 \(\delta_0\) 将 \(\mathcal E_0\) 赋成 \(\mathcal A(G)\)-余模,
从而定义 \(G\) 在 \(\mathcal E_0\) 上的作用 \(\rho_0\).
称其为 \(\rho\) 在 \(\mathcal E_0\) 上诱导的作用,
并称 \(\mathcal E_0\) 在 \(\rho\) 下稳定.}

这由定义及\hyperref[par:sga3-VIB-11.6]{11.6} 立即得到.
但须注意, 一般而言典范映射
\[
  \mathrm W(\mathcal E_0)\longrightarrow\mathrm W(\mathcal E)
\]
并非单态射.

\medskip
\phantomsection\label{rem:sga3-VIB-11.7.bis}%
\noindent\textbf{注 11.7.bis.---}%
\footnote{\label{note:sga3-VIB-111}编者注: 此注为新增内容, 推广了\SGARefLink{lem:sga3-VIB-11.7}{11.7},
并将在\SGARefLink{prop:sga3-VIB-11.10.bis}{11.10.b}is 中用到.}
设 \(G\) 为平坦 \(S\)-群, \(\mathcal E\) 为拟凝聚
\(G\)-\(\mathcal O_S\)-模.
用 \(f\) 表示态射 \(G\to S\), 用 \(\delta\) 表示态射
\[
  \mathcal E\longrightarrow f_*f^*(\mathcal E)
\]
它定义于\hyperref[par:sga3-VIB-11.6]{11.6}.
设 \(\mathcal E_0\) 为 \(\mathcal E\) 的拟凝聚 \(\mathcal O_S\)-子模.
由于 \(f\) 平坦, \(f_*f^*(\mathcal E_0)\) 是 \(f_*f^*(\mathcal E)\) 的
\(\mathcal O_S\)-子模, 对 \(\varphi=f\times f\) 也有同样结论.
因此, 若 \(\delta\) 在 \(\mathcal E_0\) 上的限制 \(\delta_0\)
经由 \(f_*f^*(\mathcal E_0)\) 分解,
则它将 \(\mathcal E_0\) 赋成 \(G\)-\(\mathcal O_S\)-模.
此时称 \(\mathcal E_0\) 为 \(\mathcal E\) 的 \(G\)-稳定子模.
